\documentclass[12pt,a4paper]{article}
\usepackage[utf8]{inputenc}
\usepackage[T5]{fontenc}
\usepackage{amsmath, amssymb}
\usepackage{graphicx}
\usepackage{ifthen}

\newboolean{showsolutions}
\setboolean{showsolutions}{true}

% --- ĐỊNH NGHĨA CÁC LỆNH ẢO CHO CÔNG CỤ LATEX2JSON CỦA WEBSITE ---
\newcommand{\doctitle}[1]{\begin{center}\Large\textbf{#1}\end{center}}
\newcommand{\docdesc}[1]{\textbf{Mô tả:} #1}
\newcommand{\docgrade}[1]{\textbf{Lớp:} #1}
\newcommand{\docstatus}[1]{\textbf{Trạng thái:} #1}
\newcommand{\doctype}[1]{\textbf{Loại:} #1}
\newcommand{\doctopics}[1]{\textbf{Chủ đề:} #1}

\newenvironment{textblock}{}{}

\newenvironment{lesson}[2]{
	\noindent\textbf{\Large #1}\\
	\textit{#2}
}{}

\newcommand{\image}[3]{
	\begin{center}
		\includegraphics[width=#1\textwidth]{#3} \\
		\textit{#2}
	\end{center}
}

\newenvironment{quiz}[1]{
	\noindent\textbf{\Large Kiểm tra: #1}
}{}

\newenvironment{mcq}[4]{
	\noindent\textbf{Câu hỏi:} #1 \\
	\ifthenelse{\boolean{showsolutions}}{
		\textit{(Đáp án đúng: #2, Điểm: #3) - Giải thích: #4}
	}{}
	\begin{itemize}
	}{
	\end{itemize}
}
\newcommand{\option}[1]{\item #1}

\newenvironment{truefalse}[3]{
	\noindent\textbf{Đúng/Sai:} #1 \\
	\ifthenelse{\boolean{showsolutions}}{
		\textit{(Điểm: #2) - Giải thích: #3}
	}{}
	\begin{itemize}
	}{
	\end{itemize}
}
\newcommand{\statement}[2]{\item [\textbf{#1}] #2}

\newcommand{\shortanswer}[4]{
    \noindent\textbf{Trả lời ngắn (Điểm: #3):}

    \noindent #1

	\ifthenelse{\boolean{showsolutions}}{
		\noindent\textit{Đáp án: #2 - Giải thích:}
		\noindent #4
	}{}
}

\newcommand{\essay}[3]{
    \noindent\textbf{Tự luận (Điểm: #2):}
    
    \noindent #1
    
	\ifthenelse{\boolean{showsolutions}}{
		\noindent\textbf{Gợi ý / Đáp án:}
		\noindent #3
	}{}
}

\begin{document}

\doctitle{BÀI 3. GÓC GIỮA HAI ĐƯỜNG THẲNG}
\docdesc{Tóm tắt lý thuyết và bài tập rèn luyện góc giữa hai đường thẳng trong không gian kèm lời giải chi tiết - Hình học 11}
\docgrade{11}
\docstatus{draft}
\doctype{normal}
\doctopics{Hinh-hoc-khong-gian, Quan-he-vuong-goc, Goc-giua-hai-duong-thang}

% =========================================================================
% PHẦN A: TÓM TẮT LÝ THUYẾT
% =========================================================================

\begin{lesson}{A. Tóm tắt lý thuyết}{Kiến thức cơ bản và các phương pháp xác định góc giữa hai đường thẳng}

\textbf{1. Cách xác định góc giữa hai đường thẳng}

- Chọn điểm $O$ tùy ý.

- Qua $O$ kẻ hai đường thẳng $a', b'$ sao cho $a' \parallel a, b' \parallel b$.

- Khi đó góc giữa hai đường thẳng $a, b$ chính bằng góc giữa hai đường thẳng $a', b'$.

\image{0.6}{}{lt_1.png}

\textbf{2. Các phương pháp tính góc}

1. Sử dụng hệ thức lượng trong tam giác:

Xét $\triangle ABC$, đặt $a = BC, b = AC, c = AB$. Ta có:

- (a) Định lí sin:
$$\frac{a}{\sin A} = \frac{b}{\sin B} = \frac{c}{\sin C}$$

- (b) Định lí cos:
$$\cos A = \frac{b^2 + c^2 - a^2}{2bc}$$

\image{0.6}{}{lt_2.png}

2. Tính theo véc-tơ chỉ phương:

Gọi $\varphi$ là góc giữa hai đường thẳng $a$ và $b$. Nếu đường thẳng $a, b$ có véc-tơ chỉ phương lần lượt là $\vec{u}_1, \vec{u}_2$ thì:
$$\cos \varphi = \frac{|\vec{u}_1 \cdot \vec{u}_2|}{|\vec{u}_1| \cdot |\vec{u}_2|}$$

*Chú ý:* $0^\circ \le \varphi \le 90^\circ$.

3. $AB \perp CD \Leftrightarrow \vec{AB} \cdot \vec{CD} = 0$.

4. Nếu $a$ và $b$ song song hoặc trùng nhau thì góc giữa chúng là $\varphi = 0^\circ$.

\end{lesson}

% =========================================================================
% PHẦN B: BÀI TẬP RÈN LUYỆN
% =========================================================================

\begin{lesson}{B. Bài tập rèn luyện}{Dạng 3.1: Tính góc giữa hai đường thẳng}

\textbf{DẠNG 3.1. Tính góc giữa hai đường thẳng}

\textbf{Phương pháp giải:}
- Cách xác định góc giữa hai đường thẳng chéo nhau $a$ và $b$: Chọn điểm $O$ thích hợp, rồi kẻ hai đường thẳng đi qua điểm $O$: $a' \parallel a$ và $b' \parallel b$. Góc giữa hai đường thẳng $a$ và $b$ là góc giữa hai đường thẳng $a'$ và $b'$.

\end{lesson}

\begin{quiz}{Bài tập rèn luyện: Dạng 3.1}

\essay{Cho hình chóp $S.ABC$ có $SA = SB = SC = AB = AC = a, BC = a\sqrt{2}$. Tính góc giữa hai đường thẳng $SC$ và $AB$.}{1}{\image{0.6}{}{lt_3.png}
- Gọi $E, F, G$ lần lượt là trung điểm của $BC, AC, SA$.
- Ta có $EF \parallel AB$ và $FG \parallel SC$.
- Do đó $(SC, AB) = (EF, FG) = \widehat{EFG}$ (hoặc bù với $\widehat{EFG}$).
- Ta có $FE = \frac{1}{2}AB = \frac{a}{2}$, $FG = \frac{1}{2}SC = \frac{a}{2}$.
- Xét $\triangle BAG$ và $\triangle CAG$:
  Ta có $AB = AC = a$, $AG$ chung, $SB = SC \Rightarrow GB = GC$ (c-g-c).
- Tam giác $GBC$ cân tại $G$ có $GE$ là trung tuyến nên đồng thời là đường cao:
  $$GE = \sqrt{BG^2 - BE^2} = \sqrt{\left(\frac{a\sqrt{3}}{2}\right)^2 - \left(\frac{a\sqrt{2}}{2}\right)^2} = \frac{a}{2}$$
  (với $BG = \frac{a\sqrt{3}}{2}$ là đường cao trong tam giác đều $SAB$).
- Xét tam giác $EFG$ có $FE = FG = GE = \frac{a}{2}$.
  Suy ra tam giác $EFG$ đều, do đó $\widehat{EFG} = 60^\circ$.
- Vậy góc giữa hai đường thẳng $SC$ và $AB$ bằng $60^\circ$.}

\essay{Cho hình hộp $ABCD.A'B'C'D'$ có độ dài tất cả các cạnh đều bằng $a > 0$ và $\widehat{BAD} = \widehat{DAA'} = \widehat{A'AB} = 60^\circ$. Gọi $M, N$ lần lượt là trung điểm của $AA', CD$. Chứng minh $MN \parallel (A'C'D)$ và tính cô-sin của góc tạo bởi hai đường thẳng $NM$ và $B'C$.}{1}{\image{0.6}{}{lt_4.png}
- Gọi $I$ là trung điểm của $DC'$.
- Trong tam giác $CDC'$, $NI$ là đường trung bình nên $NI \parallel CC'$ và $NI = \frac{1}{2}CC'$.
- Mặt khác tính chất hình hộp ta có $CC' \parallel AA'$ và $CC' = AA'$.
  Do đó $NI \parallel AA'$ và $NI = \frac{1}{2}AA' \Rightarrow NI \parallel MA'$ và $NI = MA'$.
- Tứ giác $MNIA'$ là hình bình hành nên $MN \parallel IA'$.
- Mà $IA' \subset (A'C'D)$ nên $MN \parallel (A'C'D)$.
- Vì $MN \parallel IA'$ và $CB' \parallel DA'$ nên $(CB', MN) = (DA', IA') = \widehat{DA'I}$ (hoặc $180^\circ - \widehat{DA'I}$).
- Tam giác $DAA'$ đều cạnh $a$ nên $DA' = a$.
- Xét tam giác $ABC$ có $\widehat{ABC} = 120^\circ$:
  $$AC^2 = AB^2 + BC^2 - 2AB \cdot BC \cdot \cos 120^\circ = 3a^2 \Rightarrow AC = a\sqrt{3}$$
- Tương tự trong tam giác $A'AB'$ có $AB' = a\sqrt{3}$.
- Do đó $A'C' = AC = a\sqrt{3}$ và $DC' = AB' = a\sqrt{3}$.
- Trong tam giác $A'DC'$, $A'I$ là đường trung tuyến:
  $$IA'^2 = \frac{DA'^2 + A'C'^2}{2} - \frac{DC'^2}{4} = \frac{a^2 + 3a^2}{2} - \frac{3a^2}{4} = \frac{5a^2}{4} \Rightarrow IA' = \frac{a\sqrt{5}}{2}$$
- Trong tam giác $A'DI$, có $DI = \frac{a\sqrt{3}}{2}$:
  $$\cos \widehat{DA'I} = \frac{DA'^2 + IA'^2 - DI^2}{2 \cdot DA' \cdot IA'} = \frac{a^2 + \frac{5a^2}{4} - \frac{3a^2}{4}}{2 \cdot a \cdot \frac{a\sqrt{5}}{2}} = \frac{3}{2\sqrt{5}} = \frac{3\sqrt{5}}{10} > 0$$
- Kết luận: $\cos(MN, CB') = \cos \widehat{DA'I} = \frac{3\sqrt{5}}{10}$.}

\essay{Cho hình lăng trụ tam giác đều $ABC.A'B'C'$ có $AB = 1, CC' = m$ ($m > 0$). Tìm $m$ biết rằng góc giữa hai đường thẳng $AB'$ và $BC'$ bằng $60^\circ$.}{1}{\image{0.6}{}{lt_5.png}
- Trong mặt phẳng $(ABB'A')$, kẻ $BD \parallel AB'$ ($D \in A'B'$ kéo dài).
- Khi đó $(AB', BC') = (BD, BC') = 60^\circ \Rightarrow \widehat{DBC'} = 60^\circ$ hoặc $\widehat{DBC'} = 120^\circ$.
- Vì $ABC.A'B'C'$ là lăng trụ đều nên $BB' \perp (A'B'C')$.
- Áp dụng định lý Pytago:
  $$BD = \sqrt{DB'^2 + BB'^2} = \sqrt{1 + m^2}$$
  $$BC' = \sqrt{C'B'^2 + BB'^2} = \sqrt{1 + m^2}$$
- Trong $\triangle A'B'C'$ đều cạnh $1$, $\widehat{DB'C'} = 180^\circ - 60^\circ = 120^\circ$. Áp dụng định lý cô-sin cho $\triangle DB'C'$:
  $$DC'^2 = DB'^2 + C'B'^2 - 2DB' \cdot C'B' \cdot \cos 120^\circ = 1 + 1 - 2\left(-\frac{1}{2}\right) = 3 \Rightarrow DC' = \sqrt{3}$$
- **Trường hợp 1:** Nếu $\widehat{DBC'} = 60^\circ$, tam giác $BDC'$ cân có góc $60^\circ$ nên đều:
  $$BD = DC' \Leftrightarrow \sqrt{m^2 + 1} = \sqrt{3} \Leftrightarrow m = \sqrt{2}\ (do\ m > 0)$$
- **Trường hợp 2:** Nếu $\widehat{DBC'} = 120^\circ$, áp dụng định lý cô-sin cho $\triangle BDC'$:
  $$DC'^2 = BD^2 + BC'^2 - 2BD \cdot BC' \cdot \cos 120^\circ \Rightarrow 3 = 3(1 + m^2) \Rightarrow m = 0\ (loại)$$
- Vậy $m = \sqrt{2}$.}

\essay{Cho lăng trụ $ABC.A'B'C'$ có độ dài cạnh bên bằng $2a$, đáy $ABC$ là tam giác vuông tại $A$, $AB = a, AC = a\sqrt{3}$ và hình chiếu vuông góc của đỉnh $A'$ trên mặt phẳng $(ABC)$ là trung điểm của cạnh $BC$. Tính cô-sin của góc giữa hai đường thẳng $AA', B'C$.}{1}{\image{0.6}{}{lt_6.png}
- Gọi $H$ là trung điểm $BC$, theo đề bài $A'H \perp (ABC)$.
- Mà $(ABC) \parallel (A'B'C') \Rightarrow A'H \perp (A'B'C')$.
- Tam giác $ABC$ vuông tại $A$:
  $$BC = \sqrt{AB^2 + AC^2} = 2a \Rightarrow BH = a$$
- Ta có $AA' \parallel BB'$ và $B'C \parallel BC$ nên $(AA', B'C) = (BB', BC) = \widehat{B'BH}$ (hoặc bù).
- Vì tam giác $ABC$ vuông tại $A$ nên $AH = \frac{BC}{2} = a$.
- Trong tam giác vuông $A'HA$:
  $$A'H = \sqrt{AA'^2 - AH^2} = \sqrt{4a^2 - a^2} = a\sqrt{3}$$
- Trong tam giác $A'B'H$ vuông tại $A'$ (do $A'H \perp (A'B'C')$):
  $$HB' = \sqrt{A'B'^2 + A'H^2} = \sqrt{a^2 + 3a^2} = 2a$$
- Áp dụng định lý cô-sin cho tam giác $B'BH$:
  $$\cos \widehat{B'BH} = \frac{BB'^2 + BH^2 - B'H^2}{2 \cdot BB' \cdot BH} = \frac{4a^2 + a^2 - 4a^2}{2 \cdot 2a \cdot a} = \frac{1}{4} > 0$$
- Vậy cô-sin góc giữa hai đường thẳng $AA'$ và $B'C$ là $\frac{1}{4}$.}

\essay{Cho hình chóp $S.ABCD$ có đáy $ABCD$ là hình vuông cạnh $2a$, $SA = a, SB = a\sqrt{3}$ và mặt phẳng $(SAB)$ vuông góc với mặt phẳng đáy. Gọi $M, N$ lần lượt là trung điểm của các cạnh $AB, BC$. Tính cô-sin của góc giữa hai đường thẳng $SM, DN$.}{1}{\image{0.6}{}{lt_7.png}
- Hạ $SH \perp AB$ tại $H$. Vì $(SAB) \perp (ABCD)$ nên $SH \perp (ABCD)$.
- Trong mặt phẳng $(ABCD)$, từ $M$ kẻ $ME \parallel DN$ với $E \in AD$. Góc giữa $SM$ và $DN$ là góc giữa $SM$ và $ME$.
- Xét tam giác $SAB$ có $AB^2 = 4a^2$, $SA^2 + SB^2 = a^2 + 3a^2 = 4a^2 \Rightarrow \triangle SAB$ vuông tại $S$.
- Hệ thức lượng trong $\triangle SAB$ vuông tại $S$:
  $$\frac{1}{SH^2} = \frac{1}{SA^2} + \frac{1}{SB^2} = \frac{4}{3a^2} \Rightarrow SH = \frac{a\sqrt{3}}{2}$$
- Tam giác $SHA$ vuông tại $H$:
  $$HA = \sqrt{SA^2 - SH^2} = \frac{a}{2}$$
- Gọi $K$ là trung điểm của $AD$, ta có $ME \parallel BK \parallel DN$.
  Do đó $ME$ là đường trung bình của tam giác $ABK$:
  $$AE = \frac{1}{2}AK = \frac{a}{2}, \quad ME = \frac{1}{2}BK = \frac{1}{2}\sqrt{AB^2 + AK^2} = \frac{a\sqrt{5}}{2}$$
- Tam giác $HAE$ vuông tại $A$:
  $$HE = \sqrt{AH^2 + AE^2} = \sqrt{\frac{a^2}{4} + \frac{a^2}{4}} = \frac{a\sqrt{2}}{2}$$
- Tam giác $SHE$ vuông tại $H$:
  $$SE = \sqrt{SH^2 + HE^2} = \sqrt{\frac{3a^2}{4} + \frac{2a^2}{4}} = \frac{a\sqrt{5}}{2}$$
- Lại có $SM = \frac{AB}{2} = a$. Áp dụng định lý cô-sin cho tam giác $SME$:
  $$\cos \widehat{SME} = \frac{SM^2 + ME^2 - SE^2}{2 \cdot SM \cdot ME} = \frac{a^2 + \frac{5a^2}{4} - \frac{5a^2}{4}}{2 \cdot a \cdot \frac{a\sqrt{5}}{2}} = \frac{1}{\sqrt{5}} = \frac{\sqrt{5}}{5} > 0$$
- Vậy cô-sin của góc giữa hai đường thẳng $SM$ và $DN$ bằng $\frac{\sqrt{5}}{5}$.}

\essay{Cho khối chóp $S.ABCD$ có $ABCD$ là hình vuông cạnh $a$, $SA = a\sqrt{3}$ và $SA$ vuông góc với mặt phẳng đáy. Tính cô-sin của góc giữa hai đường thẳng $SB, AC$.}{1}{\image{0.6}{}{lt_9.png}
- Qua $B$ kẻ đường thẳng $Bx \parallel AC$, $Bx$ cắt đường thẳng $AD$ tại $F$.
- Góc giữa $AC$ và $SB$ chính là góc giữa $Bx$ và $SB$.
- Tứ giác $AFBC$ là hình bình hành nên $AF = BC = a$ và $FB = AC = a\sqrt{2}$.
- Xét tam giác vuông $SAB$ và $SAF$:
  $$SB = \sqrt{SA^2 + AB^2} = 2a, \quad SF = \sqrt{SA^2 + AF^2} = 2a$$
  Do đó tam giác $SBF$ cân tại $S$.
- Hạ $SH \perp FB$ tại $H \Rightarrow H$ là trung điểm $FB$:
  $$HB = \frac{FB}{2} = \frac{a\sqrt{2}}{2}$$
- Trong tam giác vuông $SHB$:
  $$\cos \widehat{SBF} = \frac{HB}{SB} = \frac{\frac{a\sqrt{2}}{2}}{2a} = \frac{\sqrt{2}}{4} > 0$$
- Vậy cô-sin của góc giữa hai đường thẳng $SB$ và $AC$ bằng $\frac{\sqrt{2}}{4}$.}

\essay{Cho hình chóp $S.ABCD$ có tất cả các cạnh đều bằng $a$, đáy là hình vuông. Gọi $N$ là trung điểm $SB$. Tính góc giữa $AN$ và $CN$, $AN$ và $SD$.}{1}{\image{0.6}{}{lt_10.png}
- Gọi $O = AC \cap BD$. Các tam giác $SAB, SBC$ đều cạnh $a$ nên $AN = CN = \frac{a\sqrt{3}}{2}$.
- Đáy $ABCD$ là hình vuông cạnh $a \Rightarrow AC = a\sqrt{2}$.
- Áp dụng định lý cô-sin cho tam giác $ANC$:
  $$\cos \widehat{ANC} = \frac{AN^2 + CN^2 - AC^2}{2 \cdot AN \cdot CN} = \frac{\frac{3a^2}{4} + \frac{3a^2}{4} - 2a^2}{2 \cdot \frac{3a^2}{4}} = -\frac{1}{3} < 0$$
  Suy ra $(AN, CN) = 180^\circ - \widehat{ANC} = \arccos\left(\frac{1}{3}\right)$.
- Trong tam giác $BDS$, $ON$ là đường trung bình nên $ON \parallel SD \Rightarrow (AN, SD) = (AN, NO) = \widehat{ANO}$.
- Ta có $ON = \frac{a}{2}, AO = \frac{a\sqrt{2}}{2}, AN = \frac{a\sqrt{3}}{2}$. Áp dụng định lý cô-sin cho $\triangle ANO$:
  $$\cos \widehat{ANO} = \frac{AN^2 + ON^2 - AO^2}{2 \cdot AN \cdot ON} = \frac{\frac{3a^2}{4} + \frac{a^2}{4} - \frac{2a^2}{4}}{2 \cdot \frac{a\sqrt{3}}{2} \cdot \frac{a}{2}} = \frac{\sqrt{3}}{3} > 0$$
  Suy ra $\widehat{ANO} = \arccos\left(\frac{\sqrt{3}}{3}\right)$.
- Kết luận: $(AN, CN) = \arccos\left(\frac{1}{3}\right)$ và $(AN, SD) = \arccos\left(\frac{\sqrt{3}}{3}\right)$.}

\end{quiz}

\end{document}
